
\subsubsection{2.1 Multi-Dimension Trustworthiness Evaluation}

DecodingTrust \citep{Wangetal2023} evaluates GPT-3.5 and GPT-4 across eight trust dimensions including adversarial robustness, calibration, fairness, and privacy. The study evaluates dimensions independently and does not compute cross-dimension predictive correlations or residualize any metric against capability. TrustLLM \citep{Sunetal2024} extends this approach to 16 LLMs across eight dimensions, reporting dimension scores independently. Neither framework tests whether adversarial fragility predicts calibration error after capability control.

A separate line of work [ctlllll, 2024] demonstrates that standard capability benchmarks cluster into two principal components explaining 97.4\% of variance, and that two LASSO-selected benchmarks achieve 0.94 Spearman correlation with human Elo ratings. This study applies a methodologically similar PCA capability compression approach but uses it as a confound control variable toward a different research question — the cross-dimension prediction between adversarial robustness and calibration.

\subsubsection{2.2 Adversarial Robustness Benchmarking}

AdvGLUE \citep{Wangetal2021} provides the primary adversarial robustness measure used in this study. It applies 14 attack methods, including textual perturbations, synonym substitutions, and paraphrase attacks, on GLUE-derived NLU tasks. Prior work applying AdvGLUE has generally reported robustness scores in isolation without correlating them with other trust dimensions. Know Thy Judge \citep{Eirasetal2025} documents that safety-specialized models show a +0.24 false negative rate increase under style perturbations, illustrating that adversarial sensitivity extends across model types.

\subsubsection{2.3 LLM Calibration}

Guo et al. [2017] showed that modern neural networks tend to be overconfident and proposed temperature scaling as a post-hoc calibration method, finding that larger networks are more overconfident. Minderer et al. [2021] revisited calibration for large pretrained models and found improved out-of-the-box calibration relative to earlier architectures. TruthfulQA \citep{Linetal2021} documents an inverse scaling result where larger models are less truthful, motivating the use of a composite capability index rather than raw parameter count as a confound control.

Ouyang et al. [2022] and Ziegler et al. [2019] demonstrate that RLHF and instruction tuning improve model reliability on in-distribution tasks while simultaneously introducing specific adversarial failure modes. This dual effect provides the theoretical basis for the calibration–robustness trade-off interpretation discussed in Section 6.1.

\subsubsection{2.4 Research Gap}

No prior study has computed partial correlations between a capability-residualized adversarial fragility measure and calibration error across a diverse multi-family LLM set. DecodingTrust and TrustLLM provide benchmark infrastructure but do not perform predictive correlation analyses across dimensions. Calibration research identifies scale and training-regime effects on ECE but does not connect these to adversarial robustness. Adversarial robustness work validates robustness signals without correlating them to calibration. This study combines these streams through the RI construct, reporting the direction of the RI–ECE partial correlation as the central empirical question.

